% !Mode:: "TeX:UTF-8"
% pages/abstract.tex

%%%--- “中文摘要” --- %%%%%%%%%%%%%%%%%

% \pagestyle{plain}
% \newpage
% \newpage  \pagestyle{plain} \myplain
% \newgeometry{left=2.6cm, right=2.1cm, top=2.5cm, bottom=2.1cm, includefoot}
% \linespread{1.25}

\begin{flushleft}
{\zihao{3}\heiti
    \makebox[5em][s]{\textbf{论文题目：}}\textbf{\the\CTitle}\\
    \makebox[5em][s]{\textbf{专业名称：}}\textbf{\the\Cmajor}\\
    \makebox[5em][s]{\textbf{研究生：}}\textbf{\the\StudentAuthor}\hfill \makebox[3.5em][s]{\textbf{签名：}}\thickuline{\quad\quad\quad\quad\quad\quad}\\
    \makebox[5em][s]{\textbf{指导老师：}}\textbf{\the\Csupervisor~~\the\CsupervisorTitle}\hfill \makebox[3.5em][s]{\textbf{签名：}}\thickuline{\quad\quad\quad\quad\quad\quad}\\
}
\end{flushleft}

\begin{cnabstract}
博士学位论文摘要是对论文研究内容、方法与主要成果的高度概括，是读者了解论文核心贡献的重要部分。摘要应能够在不阅读全文的情况下，使读者清晰了解研究问题、研究方法及主要结论。一般而言，博士论文摘要应包括以下几个方面的内容。研究背景与问题提出：简要介绍研究领域的发展背景及存在的关键问题，说明本研究开展的必要性与研究意义。该部分应突出研究问题，而不宜进行过多的文献综述；研究目标与研究内容：概述论文的主要研究目标，并简要说明论文围绕该目标开展的核心研究内容或研究思路；研究方法与技术路线：介绍论文所采用的主要研究方法、理论模型或技术框架，使读者能够理解研究工作的基本思路与方法特点；主要研究成果与创新点：重点概括论文取得的主要研究成果，包括提出的新方法、新模型或新理论，并突出论文的创新性贡献；结论与应用价值：简要总结研究结论，并说明研究成果的理论意义或实际应用价值。

在写作形式上，摘要应具有独立性和完整性，避免使用“本文”“本章”等指代性较强的表达，同时一般不引用文献、不出现图表和公式。摘要内容应逻辑清晰、语言精炼，通常控制在 800–1200 字左右（中文摘要）。摘要之后应给出 3–5 个关键词（Keywords），用于反映论文研究的核心主题。关键词应尽量选用学科领域内通用、规范的术语，并按照重要程度进行排列。

\par\vspace{\baselineskip}
\cnkeywords{博士学位论文；论文模板；格式规范；图表使用；附件引入}
\end{cnabstract}

\cleardoublepage % 结束当前页
% \restoregeometry % 恢复之前的页边距设置

%%%--- “英文摘要” --- %%%%%%%%%%%%%%%%%

\begin{flushleft}
{
\renewcommand{\baselinestretch}{1}
\zihao{4}\textsf{\textbf{Title: \the\ETitle}\\
    \vspace*{14pt}
    \textbf{Major: \the\Emajor}\\
    \vspace*{14pt}
    \textbf{Name: \the\EStudentAuthor}\hfill \textbf{Signature:} \thickuline{\quad\quad\quad\quad\quad\quad}\\
    \vspace*{14pt}
    \textbf{Supervisor: \the\EsupervisorTitle~~\the\Esupervisor}\hfill \textbf{Signature:} \thickuline{\quad\quad\quad\quad\quad\quad}\\
}}
\end{flushleft}

\begin{enabstract}

The doctoral dissertation abstract provides a concise summary of the thesis's research content, methodology, and primary findings, serving as a crucial component for readers to grasp the core contributions of the work. It should enable readers to clearly understand the research question, methodology, and main conclusions without reading the full text. Generally, a doctoral dissertation abstract should include the following elements. Research Background and Problem Statement: Briefly introduce the developmental context of the research field and the key existing problems, explaining the necessity and significance of this study. This section should highlight the research question without excessive literature review; Research Objectives and Content: Outline the primary research objectives and briefly describe the core research content or approach developed around these objectives; Research Methods and Technical Approach: Introduce the primary research methods, theoretical models, or technical frameworks employed, enabling readers to grasp the fundamental methodology and methodological characteristics of the work. Key Findings and Innovations: Summarize the main research outcomes, including novel methods, models, or theories proposed, emphasizing the innovative contributions of the thesis. Conclusions and Practical Value: Briefly summarize the research conclusions and articulate the theoretical significance or practical application value of the findings.

In terms of writing format, the abstract should be self-contained and complete, avoiding highly referential expressions such as “this paper” or “this chapter.” Generally, it should not cite references, include figures or tables, or present mathematical formulas. The abstract should be logically coherent and concisely written, typically limited to 800–1200 words (for Chinese abstracts). Following the abstract, provide 3–5 keywords (Keywords) reflecting the core themes of the paper's research. Keywords should prioritize standardized terminology within the discipline and be listed in order of importance.

\par\vspace{\baselineskip}
\enkeywords{Doctoral Dissertation; Dissertation Template; Formatting Guidelines; Use of Figures and Tables; Inclusion of Appendices}
\end{enabstract}

\cleardoublepage % 结束当前页
% \restoregeometry % 恢复之前的页边距设置